\documentclass[10pt,a4paper]{amsart}
\usepackage{amsmath,amssymb,amsthm,graphicx}
% For commutative diagrams we use Xy-pic
\usepackage[all]{xy}

% We use 2cell for 2-commutative diagrams.
\xyoption{2cell}
% Theorem environments.
%
\theoremstyle{plain}
\newtheorem{theorem}[subsection]{Theorem}
\newtheorem{proposition}[subsection]{Proposition}
\newtheorem{lemma}[subsection]{Lemma}

\theoremstyle{definition}
\newtheorem{definition}[subsection]{Definition}
\newtheorem{example}[subsection]{Example}
\newtheorem{exercise}[subsection]{Exercise}
\newtheorem{situation}[subsection]{Situation}

\theoremstyle{remark}
\newtheorem{remark}[subsection]{Remark}
\newtheorem{remarks}[subsection]{Remarks}

\numberwithin{equation}{subsection}

% Macros
%
\def\lim{\mathop{\mathrm{lim}}\nolimits}
\def\colim{\mathop{\mathrm{colim}}\nolimits}
\def\Spec{\mathop{\mathrm{Spec}}}
\def\Hom{\mathop{\mathrm{Hom}}\nolimits}
\def\Ext{\mathop{\mathrm{Ext}}\nolimits}
\def\SheafHom{\mathop{\mathcal{H}\!\mathit{om}}\nolimits}
\def\SheafExt{\mathop{\mathcal{E}\!\mathit{xt}}\nolimits}
\def\Sch{\mathit{Sch}}
\def\Mor{\mathop{\mathrm{Mor}}\nolimits}
\def\Ob{\mathop{\mathrm{Ob}}\nolimits}
\def\Sh{\mathop{\mathit{Sh}}\nolimits}
\def\NL{\mathop{N\!L}\nolimits}
\def\CH{\mathop{\mathrm{CH}}\nolimits}
\def\proetale{{pro\text{-}\acute{e}tale}}
\def\etale{{\acute{e}tale}}
\def\QCoh{\mathit{QCoh}}
\def\Ker{\mathop{\mathrm{Ker}}}
\def\Im{\mathop{\mathrm{Im}}}
\def\Coker{\mathop{\mathrm{Coker}}}
\def\Coim{\mathop{\mathrm{Coim}}}

% Boxtimes
%
\DeclareMathSymbol{\boxtimes}{\mathbin}{AMSa}{"02}

%
% Macros for moduli stacks/spaces
%
\def\QCohstack{\mathcal{QC}\!\mathit{oh}}
\def\Cohstack{\mathcal{C}\!\mathit{oh}}
\def\Spacesstack{\mathcal{S}\!\mathit{paces}}
\def\Quotfunctor{\mathrm{Quot}}
\def\Hilbfunctor{\mathrm{Hilb}}
\def\Curvesstack{\mathcal{C}\!\mathit{urves}}
\def\Polarizedstack{\mathcal{P}\!\mathit{olarized}}
\def\Complexesstack{\mathcal{C}\!\mathit{omplexes}}
% \Pic is the operator that assigns to X its picard group, usage \Pic(X)
% \Picardstack_{X/B} denotes the Picard stack of X over B
% \Picardfunctor_{X/B} denotes the Picard functor of X over B
\def\Pic{\mathop{\mathrm{Pic}}\nolimits}
\def\Picardstack{\mathcal{P}\!\mathit{ic}}
\def\Picardfunctor{\mathrm{Pic}}
\def\Deformationcategory{\mathcal{D}\!\mathit{ef}}
\usepackage{cleveref}
\crefname{originalreference}{original reference}{original references}
\Crefname{originalreference}{Original reference}{Original references}
\makeatletter
\newlabel{algebra-example-localize-at-prime}{{02C5}{source}{}{sourcepointer-1}{}}
\newlabel{algebra-example-localize-at-prime@cref}{{[originalreference][0][]02C5}{[1][0][]source}{}{}{}}
\makeatother
\crefname{originalsection}{original section}{original sections}
\Crefname{originalsection}{Original section}{Original sections}
\makeatletter
\newlabel{more-algebra-section-determinants-finite-length}{{0GSX}{source}{}{sourcepointer-2}{}}
\newlabel{more-algebra-section-determinants-finite-length@cref}{{[originalsection][0][]0GSX}{[1][0][]source}{}{}{}}
\makeatother
\crefname{originallemma}{original lemma}{original lemmas}
\Crefname{originallemma}{Original lemma}{Original lemmas}
\makeatletter
\newlabel{more-algebra-lemma-multiplication}{{0GSZ}{source}{}{sourcepointer-3}{}}
\newlabel{more-algebra-lemma-multiplication@cref}{{[originallemma][0][]0GSZ}{[1][0][]source}{}{}{}}
\makeatother
\crefname{originallemma}{original lemma}{original lemmas}
\Crefname{originallemma}{Original lemma}{Original lemmas}
\makeatletter
\newlabel{more-algebra-lemma-flat-base-change-det}{{0GT0}{source}{}{sourcepointer-4}{}}
\newlabel{more-algebra-lemma-flat-base-change-det@cref}{{[originallemma][0][]0GT0}{[1][0][]source}{}{}{}}
\makeatother
\crefname{originallemma}{original lemma}{original lemmas}
\Crefname{originallemma}{Original lemma}{Original lemmas}
\makeatletter
\newlabel{more-algebra-lemma-ses}{{0GSY}{source}{}{sourcepointer-5}{}}
\newlabel{more-algebra-lemma-ses@cref}{{[originallemma][0][]0GSY}{[1][0][]source}{}{}{}}
\makeatother
\title{Base change of the tame symbol along a flat local homomorphism of 1-dimensional Noetherian local rings}
\author{The Stacks Project Authors}
\date{stacks-project:chow.tex:local-snapshot-20261011; GNU FDL 1.2 or later}
\begin{document}
\maketitle
\noindent\textit{Editorial note.} Counted proposition: Tag 0EPG, chow.tex lines 999-1008 with its own proof at lines 1010-1086 (statement immediately followed by its proof; no non-adjacent pairing).\par\medskip
\noindent\textit{Editorial note.} The context blocks reproduce the author's definition of the tame symbol with its defining formula (Tag 0EAQ) and the statement of the factorization lemma Tag 0EAG; the latter is also counted as a separate task of this batch, so its proof is not duplicated here.\par\medskip
\noindent\textit{Editorial note.} The displayed identity labelled equation-to-prove at chow.tex lines 1055-1061 is created inside the proof and is not an external reference.\par\medskip
\noindent\textit{Editorial note.} External references are navigation pointers to stable Stacks tags; the determinant computation that uses them is retained in full in the counted proof.\par\medskip
\noindent\textit{Allowed standard prerequisites (editorial).}\begin{itemize}
\item The tame symbol of a 1-dimensional Noetherian local ring. Stacks tag 0EAQ (retained context, chow.tex lines 759-827). Fixes the objects partial\_A(a\_1, a\_2) in kappa(m) and the local-factorization recipe that the statement compares with partial\_\{A'\}(a\_1, a\_2).
\item Local factorization after a finite extension. Stacks tag 0EAG (retained context, chow.tex lines 641-655; counted as a separate task of this batch). Supplies the finite extension A contained in B and the factorizations a\_i = u\_\{i,j\} pi\_j\textasciicircum{}\{e\_\{i,j\}\} at the maximal ideals of B, used by the author at lines 1013-1020.
\item Determinants of endomorphisms of finite-length modules, their multiplicativity and their flat base change. More on Algebra, Section 0GSX and Lemmas 0GSZ, 0GT0, 0GSY. Rewrites Norm(u)\textasciicircum{}(m\_j) as det\_kappa(u : M to M) for M = B\_j/pi\_j B\_j, moves it to det\_\{kappa'\}(u : M tensor A' to M tensor A') and splits it over the maximal ideals of B' = A' tensor\_A B (author's lines 1062-1085).
\item Identification of the residue field extension through the finite flat extension. Stacks tag 02M0 (Algebra, Lemma 02M0), used in the retained context. Underlies the length and degree bookkeeping m\_j, m'\_\{j,l\} that the comparison of the two displayed formulae relies on.
\end{itemize}\medskip
\ifdefined\hypertarget\hypertarget{sourcepointer-1}{}\fi
\noindent\textit{Source reference.} Original reference 02C5: sources/stacks-chow/tags.txt, source lines 3072--3072. The total ring of fractions Q(A) is introduced at chow.tex lines 763-765 by a pointer to this example; the counted proof never uses Q(A), only nonzerodivisors of A.\par\medskip
\ifdefined\hypertarget\hypertarget{sourcepointer-2}{}\fi
\noindent\textit{Source reference.} Original section 0GSX: sources/stacks-chow/tags.txt, source lines 20947--20947. Navigation pointer to the construction of determinants of endomorphisms of finite-length modules, invoked at chow.tex line 1064; the author's determinant comparison itself is retained in full.\par\medskip
\ifdefined\hypertarget\hypertarget{sourcepointer-3}{}\fi
\noindent\textit{Source reference.} Original lemma 0GSZ: sources/stacks-chow/tags.txt, source lines 20949--20949. Identifies the norm power with a determinant, cited at chow.tex lines 1066 and 1083; a standard input to the retained computation, not a substitute for it.\par\medskip
\ifdefined\hypertarget\hypertarget{sourcepointer-4}{}\fi
\noindent\textit{Source reference.} Original lemma 0GT0: sources/stacks-chow/tags.txt, source lines 20950--20950. Flat base change of the determinant construction, cited at chow.tex line 1072; the flatness hypothesis and its use are stated by the author, and the comparison remains the counted content.\par\medskip
\ifdefined\hypertarget\hypertarget{sourcepointer-5}{}\fi
\noindent\textit{Source reference.} Original lemma 0GSY: sources/stacks-chow/tags.txt, source lines 20948--20948. Multiplicativity of the determinant construction in short exact sequences, cited at chow.tex line 1085 to conclude the displayed identity.\par\medskip
\begin{lemma}
\label{lemma-prepare-tame-symbol}
Let $(A, \mathfrak m)$ be a Noetherian local ring of dimension $1$.
Let $r \geq 2$ and let $a_1, \ldots, a_r \in A$ be nonzerodivisors
not all units.
Then there exist
\begin{enumerate}
\item a finite ring extension $A \subset B$ with
$B/A$ annihilated by a power of $\mathfrak m$,
\item for each maximal ideal $\mathfrak m_j \subset B$
a nonzerodivisor $\pi_j \in B_j = B_{\mathfrak m_j}$, and
\item factorizations $a_i = u_{i, j} \pi_j^{e_{i, j}}$ in $B_j$
with $u_{i, j} \in B_j$ units and $e_{i, j} \geq 0$.
\end{enumerate}
\end{lemma}

\section{Tame symbols}
\label{section-tame-symbol}

\noindent
Consider a Noetherian local ring $(A, \mathfrak m)$ of dimension $1$.
We denote $Q(A)$ the total ring of fractions of $A$, see
Algebra, Example \ref{algebra-example-localize-at-prime}.
The {\it tame symbol} will be a map
$$
\partial_A(-, -) : Q(A)^* \times Q(A)^* \longrightarrow \kappa(\mathfrak m)^*
$$
satisfying the following properties:
\begin{enumerate}
\item $\partial_A(f, gh) = \partial_A(f, g) \partial_A(f, h)$
\label{item-bilinear}
for $f, g, h \in Q(A)^*$,
\item $\partial_A(f, g) \partial_A(g, f) = 1$
\label{item-skew}
for $f, g \in Q(A)^*$,
\item $\partial_A(f, 1 - f) = 1$
\label{item-1-x}
for $f \in Q(A)^*$ such that $1 - f \in Q(A)^*$,
\item $\partial_A(aa', b) = \partial_A(a, b)\partial_A(a', b)$
\label{item-bilinear-better}
and $\partial_A(a, bb') = \partial_A(a, b)\partial_A(a, b')$
for $a, a', b, b' \in A$ nonzerodivisors,
\item $\partial_A(b, b) = (-1)^m$
\label{item-skew-better}
with $m = \text{length}_A(A/bA)$
for $b \in A$ a nonzerodivisor,
\item $\partial_A(u, b) = u^m \bmod \mathfrak m$
\label{item-normalization}
with $m = \text{length}_A(A/bA)$ for $u \in A$ a unit and
$b \in A$ a nonzerodivisor, and
\item
\label{item-1-x-better}
$\partial_A(a, b - a)\partial_A(b, b) = \partial_A(b, b - a)\partial_A(a, b)$
for $a, b \in A$ such that $a, b, b - a$ are nonzerodivisors.
\end{enumerate}
Since it is easier to work with elements of $A$ we will
often think of $\partial_A$ as a map defined on pairs of
nonzerodivisors of $A$ satisfying (\ref{item-bilinear-better}),
(\ref{item-skew-better}), (\ref{item-normalization}),
(\ref{item-1-x-better}). It is an exercise to see that
setting
$$
\partial_A(\frac{a}{b}, \frac{c}{d}) =
\partial_A(a, c) \partial_A(a, d)^{-1} \partial_A(b, c)^{-1} \partial_A(b, d)
$$
we get a well defined map $Q(A)^* \times Q(A)^* \to \kappa(\mathfrak m)^*$
satisfying (\ref{item-bilinear}), (\ref{item-skew}), (\ref{item-1-x})
as well as the other properties.

\medskip\noindent
We do not claim there is a unique map with these properties.
Instead, we will give a recipe for constructing such a map.
Namely, given $a_1, a_2 \in A$ nonzerodivisors, we choose
a ring extension $A \subset B$ and local factorizations
as in Lemma \ref{lemma-prepare-tame-symbol}.
Then we define
\begin{equation}
\label{equation-tame-symbol}
\partial_A(a_1, a_2) = \prod\nolimits_j
\text{Norm}_{\kappa(\mathfrak m_j)/\kappa(\mathfrak m)}
((-1)^{e_{1, j}e_{2, j}}u_{1, j}^{e_{2, j}}u_{2, j}^{-e_{1, j}}
\bmod \mathfrak m_j)^{m_j}
\end{equation}
where $m_j = \text{length}_{B_j}(B_j/\pi_j B_j)$ and the product
is taken over the maximal ideals $\mathfrak m_1, \ldots, \mathfrak m_r$ of $B$.

\begin{lemma}
\label{lemma-tame-symbol-formally-smooth}
Let $(A, \mathfrak m, \kappa) \to (A', \mathfrak m', \kappa')$
be a local homomorphism of Noetherian local rings. Assume $A \to A'$
is flat and $\dim(A) = \dim(A') = 1$. Set
$m = \text{length}_{A'}(A'/\mathfrak mA')$.
For $a_1, a_2 \in A$ nonzerodivisors
$\partial_A(a_1, a_2)^m$ maps to $\partial_{A'}(a_1, a_2)$
via $\kappa \to \kappa'$.
\end{lemma}

\begin{proof}
If $a_1, a_2$ are both units, then $\partial_A(a_1, a_2) = 1$
and $\partial_{A'}(a_1, a_2) = 1$ and the result is true.
If not, then we can choose a ring extension $A \subset B$ and
local factorizations as in Lemma \ref{lemma-prepare-tame-symbol}.
Denote $\mathfrak m_1, \ldots, \mathfrak m_m$
be the maximal ideals of $B$. Let $\mathfrak m_1, \ldots, \mathfrak m_m$
be the maximal ideals of $B$ with residue fields $\kappa_1, \ldots, \kappa_m$.
For each $j \in \{1, \ldots, m\}$ denote $\pi_j \in B_j = B_{\mathfrak m_j}$
a nonzerodivisor such that we have factorizations
$a_i = u_{i, j}\pi_j^{e_{i, j}}$ as in the lemma.
By definition we have
$$
\partial_A(a_1, a_2) = \prod\nolimits_j
\text{Norm}_{\kappa_j/\kappa}
((-1)^{e_{1, j}e_{2, j}}u_{1, j}^{e_{2, j}}u_{2, j}^{-e_{1, j}}
\bmod \mathfrak m_j)^{m_j}
$$
where $m_j = \text{length}_{B_j}(B_j/\pi_j B_j)$.

\medskip\noindent
Set $B' = A' \otimes_A B$. Since $A'$ is flat over $A$ we see
that $A' \subset B'$ is a ring extension with $B'/A'$ annihilated
by a power of $\mathfrak m'$. Let
$$
\mathfrak m'_{j, l},\quad l = 1, \ldots, n_j
$$
be the maximal ideals of $B'$ lying over $\mathfrak m_j$. Denote
$\kappa'_{j, l}$ the residue field of $\mathfrak m'_{j, l}$. Denote
$B'_{j, l}$ the localization of $B'$ at $\mathfrak m'_{j, l}$.
As factorizations of $a_1$ and $a_2$ in $B'_{j, l}$
we use the image of the factorizations
$a_i = u_{i, j} \pi_j^{e_{i, j}}$ given to us in $B_j$.
By definition we have
$$
\partial_{A'}(a_1, a_2) = \prod\nolimits_{j, l}
\text{Norm}_{\kappa'_{j, l}/\kappa'}
((-1)^{e_{1, j}e_{2, j}}u_{1, j}^{e_{2, j}}u_{2, j}^{-e_{1, j}}
\bmod \mathfrak m'_{j, l})^{m'_{j, l}}
$$
where $m'_{j, l} = \text{length}_{B'_{j, l}}(B'_{j, l}/\pi_j B'_{j, l})$.

\medskip\noindent
Comparing the formulae above we see that it suffices to show that
for each $j$ and for any unit $u \in B_j$ we have
\begin{equation}
\label{equation-to-prove}
\left(\text{Norm}_{\kappa_j/\kappa}(u \bmod \mathfrak m_j)^{m_j}\right)^m
=
\prod\nolimits_l
\text{Norm}_{\kappa'_{j, l}/\kappa'}(u \bmod \mathfrak m'_{j, l})^{m'_{j, l}}
\end{equation}
in $\kappa'$. We are going to use the construction of determinants of
endomorphisms of finite length modules in
More on Algebra, Section \ref{more-algebra-section-determinants-finite-length}
to prove this. Set $M = B_j/\pi_j B_j$. By
More on Algebra, Lemma \ref{more-algebra-lemma-multiplication} we have
$$
\text{Norm}_{\kappa_j/\kappa}(u \bmod \mathfrak m_j)^{m_j} =
\det\nolimits_\kappa(u : M \to M)
$$
Thus, by
More on Algebra, Lemma \ref{more-algebra-lemma-flat-base-change-det},
the left hand side of (\ref{equation-to-prove}) is equal to
$\det_{\kappa'}(u : M \otimes_A A' \to M \otimes_A A')$.
We have an isomorphism
$$
M \otimes_A A' = (B_j/\pi_j B_j) \otimes_A A' =
\bigoplus\nolimits_l B'_{j, l}/\pi_j B'_{j, l}
$$
of $A'$-modules. Setting $M'_l = B'_{j, l}/\pi_j B'_{j, l}$ we see that
$\text{Norm}_{\kappa'_{j, l}/\kappa'}(u \bmod \mathfrak m'_{j, l})^{m'_{j, l}}
= \det_{\kappa'}(u_j : M'_l \to M'_l)$ by
More on Algebra, Lemma \ref{more-algebra-lemma-multiplication} again.
Hence (\ref{equation-to-prove}) holds by multiplicativity of the determinant
construction, see More on Algebra, Lemma \ref{more-algebra-lemma-ses}.
\end{proof}

\end{document}
